\documentclass[10pt,a4paper]{amsart}
\usepackage[margin=19mm]{geometry}
\usepackage{amsmath,amssymb,mathtools}
\usepackage[scr=rsfs,cal=euler]{mathalfa}
\usepackage{xfrac}
\usepackage[unicode,hidelinks,bookmarks=false]{hyperref}
\newcommand{\ie}{i.e.\ }
\newcommand{\cf}{cf.\ }
\newcommand{\resp}{respectively\ }
\newcommand{\Subsection}{\subsection}
\providecommand{\nobreakdash}{\nobreak\hbox{-}}
\setlength{\emergencystretch}{2em}
\multlinegap0pt
\usepackage{bm}
\usepackage{bbm}
\usepackage{dsfont}
\usepackage{xspace}
\usepackage{cancel}
\usepackage{mathtools}
\usepackage{amsthm}
\makeatletter
\@ifundefined{lemma}{\newtheorem{lemma}{Lemma}}{}
\@ifundefined{theorem}{\newtheorem{theorem}{Theorem}}{}
\@ifundefined{proposition}{\newtheorem{proposition}[theorem]{Proposition}}{}
\makeatother
\newcommand{\N}{\mathbb{N}}
\newcommand{\R}{\mathbb{R}}
\newcommand{\cF}{\mathcal{F}}
\newcommand{\cH}{\mathcal{H}}
\newcommand{\cK}{\mathcal{K}}
\newcommand{\eps}{\varepsilon}
\newcommand{\mP}{\mathbb{P}}
\title[Genericity of Lyapunov spectrum of bounded random compact op]{Genericity of Lyapunov spectrum of bounded random compact operators on infinite-dimensional Hilbert spaces}
\author{Thai Son Doan}
\date{Source: arXiv:2303.14359v2 (CC BY 4.0)}
\begin{document}
\maketitle

\noindent\emph{Source and licence.} Genericity of Lyapunov spectrum of bounded random compact operators on infinite-dimensional Hilbert spaces. Version arXiv:2303.14359v2 (CC BY 4.0), distributed under the Creative Commons Attribution 4.0 International licence, as recorded in the arXiv licence metadata for this exact version and shown on the abstract page https://arxiv.org/abs/2303.14359v2. Full rights notice: licenses/arxiv-2303.14359-CC-BY-4.0.txt.\par\smallskip

\noindent\textit{Editorial scope.} Counted unit: the Theorem whose source label is \texttt{HR\_Lemma} in the article (source0.tex lines 2340--2342, source label \texttt{HR\_Lemma}), delivered together with the article's own complete original proof of it (source0.tex lines 1126--1463, one \texttt{proof} environment of 338 lines). No mathematics is written, repaired or re-derived: statement and proof are the article's own text line for line. The proof is detached from the statement in the source: the article states the result at lines 2340--2342 and prints its proof at lines 1126--1463, and the proof environment's own header names the statement, so the association is the article's own. Required context retained uncounted: TechnicalProposition\_1 (lines 791--814); Millionschikov (lines 1022--1030); Goodgrowthrate (lines 1024--1029); Lowerestimate (lines 1024--1029); Perturbationalongashortperiod (lines 1035--1046). Every definition, display and label of those lines is retained unchanged. Omitted from this package and not redistributed: 1--790, 815--1021, 1030--1034, 1047--1125, 1464--2339, 2343--2667. Nothing inside a retained excerpt is omitted or altered: every retained body is delivered exactly as the article's lines are written, so no omission receipt of any kind accompanies this package. Citations: the retained excerpts carry no citation command at all, so no bibliography entry of the article is transcribed and no reference is redistributed. Reference adaptation: none was needed and none was made; every reference inside the delivered unit resolves inside it, so no unresolved reference or citation is produced. Printed slips kept literal: no printed word, symbol, hypothesis or number of the counted statement or of its proof was altered. Layout and class: the article's own class is \texttt{sn-jnl} with packages including graphicx, multirow, amsmath, amssymb, amsfonts, amsthm, mathrsfs, appendix, xcolor, textcomp, manyfoot, booktabs, algorithm, algorithmicx; the wrapper is the lane's pinned \texttt{amsart} editorial wrapper at 10pt on A4, which restates the theorem-like environments the retained text needs and adds the article's own macro definitions that the retained text needs (\texttt{\textbackslash N}, \texttt{\textbackslash R}, \texttt{\textbackslash cF}, \texttt{\textbackslash cH}, \texttt{\textbackslash cK}, \texttt{\textbackslash eps}, \texttt{\textbackslash mP}), with extra wrapper packages (bm, bbm, dsfont, xspace, cancel, mathtools). The delivered numbering is the unit's own. Assets: no retained span contains a figure environment, a graphics-inclusion command, a PDF-page inclusion, a table or a TikZ picture; no image file is copied into this package, the package declares no asset dependency and no image file is redistributed with it. Licence: version arXiv:2303.14359v2 of this article is distributed under the Creative Commons Attribution 4.0 International licence, as recorded in the arXiv licence metadata for this exact version and shown on the abstract page https://arxiv.org/abs/2303.14359v2. A full rights notice is delivered at licenses/arxiv-2303.14359-CC-BY-4.0.txt. Characters: every retained excerpt is pure ASCII, so the delivered engine, XeLaTeX with its default Latin Modern text font, reports no missing character.
\begin{proposition}\label{TechnicalProposition_1} Let $T\in\mathcal N$ and $\eps>0$ be arbitrary. Then  there exists an integer $N^*(T,\eps)\in\N$ (depending only on $T,\eps$) such that for all $N\geq N^*(T,\eps)$ there exists  a set $V\in\mathcal F$ of positive $\mP$-probability such that the sets $V,\theta V,\dots, \theta^N V$ are disjoint and the following statements hold

\noindent
(i) For any random unit vector $e:V\cup \theta^N V\rightarrow \mathcal H$, there exists $S\in\mathcal K_{\infty}(\Omega,\mathcal H)$ and a scalar random variable $\beta:\bigcup_{i=0}^{N-1}\theta^i V\rightarrow\R$ with $|\beta(\omega)|\geq \frac{\eps}{3}$ for $\omega\in \bigcup_{i=0}^{N-1}\theta^i V$ such that 
\begin{itemize}
\item [(p1)] \emph{Small perturbation}:
\[
S(\omega)=T(\omega)\hbox{ for } \omega\not\in \bigcup_{i=0}^{N-1}\theta^i V
\hbox{ and } \|S(\omega)- T(\omega)\|\leq \eps\hbox{ for } \omega\in \bigcup_{i=0}^{N-1}\theta^i V.	
\]
%\item [(p2)] \emph{Collinearity}: 
%$S^N_{\omega} h(\omega)$ is collinear to $ h(\theta^N\omega)$ %for all $\mP$-a.e. $\omega \in V$.
\item [(p2)] \emph{Existence of invariant one-dimensional random subspace along the sets $V,\theta V,\dots, \theta^N V$}: There exists an extension of $e$ to a random unit vector $e:\bigcup_{i=0}^N\theta^i V\rightarrow \cH$ such that
\[
S(\theta^i\omega)e(\theta^i\omega)=\beta(\theta^i\omega)e(\theta^{i+1}\omega)\quad \hbox{ for } \omega\in V, i=0,1,\dots,N-1.
\]
\end{itemize}

\noindent 
(ii) Let  $h:V\cup\theta^N V\rightarrow \cH$ be a random unit vector. Then, there exist $R \in \cK_{\infty}(\Omega,\cH)$, an extension of $h$ to a random unit vector $h:\Omega\rightarrow \cH$ and a scalar random variable $\gamma:\Omega\rightarrow \mathbb R$ with $|\gamma(\omega)|\geq \frac{\eps}{3}$ for $\omega\in\Omega$ such that $\|R-T\|_{\infty}\leq \eps$ and $R(\omega) h(\omega)=\gamma(\omega)h(\theta\omega)$.



\end{proposition}

\begin{proposition}\label{Millionschikov}
Let $T\in\mathcal N$ be arbitrary and let $V,\theta V,\dots, \theta^N V$ be pairwise disjoint sets satisfying assertions of Proposition \ref{TechnicalProposition_1}. Then, for any $\eps\in (0,\frac{1}{6})$ and $k\in \N$ there exist $S\in\cK_{\infty}(\Omega,\cH)$ and an invariant unit random vector $e:\Omega\rightarrow \cH$ such that $\|S-T\|_{\infty}\leq 3\eps (1+ \|T\|_{\infty})$ and
\begin{align}
\|S(\omega)e(\omega)\|
& \geq \frac{\eps^{N+1}}{(1+2\|T\|_{\infty})^{N-1}}\qquad\hbox{ for } \omega\in\Omega,\label{Lowerestimate}\\[1ex]
 \|S^{kN}_{\omega}e(\omega)\|
 & \geq  \frac{\eps^{k+4N}}{(1+2\|T\|_{\infty})^{4N}} \|S_{\omega}^{kN}\|\qquad\hbox{ for } \omega\in\Omega.\label{Goodgrowthrate}
\end{align}
\end{proposition}

\begin{lemma}\label{Perturbationalongashortperiod}
Let $T\in\cK_{\infty}(\Omega,\cH)$ be arbitrary. Let $V\subset \Omega$ be a measurable set such that $V,\theta V,\dots, \theta^{N-1} V$ are disjoint, where $N\in \mathbb N$.  Let $g: V\rightarrow \cH$ be a unit random vector. Then, for any $\eps\in (0,\frac{1}{6})$ there exists a strongly measurable map $S:\bigcup_{i=0}^{N-1} \theta^i V\rightarrow \mathcal K(\mathcal H)$ satisfying the following requirements:
\begin{itemize}
\item [(r1)] For $\omega\in \bigcup_{i=0}^{N-1} \theta^i V$, we have $\|S(\omega)-T(\omega)\|\leq 3\eps + 3\eps \|T\|_{\infty}$, 

\item [(r2)] For $\omega\in V$, we have 
 
\begin{equation}\label{Smallgrowth}
 \|S_{\omega}^N g(\omega)\|\geq \eps \|S_{\omega}^N\|\quad\hbox{and}  \quad \|S_{\omega}^N\|\geq \eps^{N+1}.
\end{equation}
\end{itemize}
\end{lemma}

\begin{theorem}\label{HR_Lemma}
Assume that $\theta$ is an ergodic invertible transformation of the non-atomic Lebesgue space $(\Omega,\mathcal F,\mathbb P)$. Then for any $U\in\cF$ with $\mP(U)>0$ and any $n > 0$ there exists a measurable set $V\subset U$ such that $\mP(V)>0$ and the sets $\theta^{-n}V,\dots,\theta^{-1}V,V,\theta V,\dots,\theta^n V$ are pairwise disjoint.
\end{theorem}

\begin{proof}[Proof of Proposition \ref{Millionschikov}]
By virtue of Halmos-Rokhlin's lemma (see Theorem \ref{HR_Lemma} in the Appendix), there exists a measurable subset $U$ of $V$ such that  $\mP(U)>0$ and the sets $U,\theta U,\dots,\theta^{(k+1)N}U$ are disjoint. By the Poincar\'{e} recurrence theorem, the set $\theta^NU$ is decomposed as the union of the following disjoint subsets $\bigcup_{i={kN+1}}^{\infty} E_i$ defined as 
\begin{equation*}\label{DecomposeofU}
E_i:=\Big\{\omega\in\theta^N U: \theta^j\omega\not\in U \hbox{ for } j=1,\dots,i-1 \hbox{ and } \theta^i\omega\in U \Big\}.    
\end{equation*}
Therefore, all the subsets $E_i,\theta E_i,\dots,\theta^i E_i$, where $i=kN+1,\dots$, are pairwise disjoint and the whose space $\Omega$ is represented as the union of the following disjoint subsets 
\begin{align}
\Omega
=&  \bigcup_{j=1}^{N-1}\theta^j U\;\cup\; \bigcup_{i=kN+1}^{\infty}  \Big(E_i\cup \theta E_i\dots\cup\theta^i E_i\Big)\label{PresentationOmega_I}\\
=&
 \bigcup_{j=0}^{N-1}\theta^j U\;\cup\; \bigcup_{i=kN+1}^{\infty}  \Big(E_i\cup \theta E_i\dots\cup\theta^{i-1} E_i\Big).\label{PresentationOmega_II}
\end{align}
The proof is divided into the construction of 
$S(\omega)$ and $e(\omega)$ and the estimate $\|S(\omega)e(\omega)\|, \|S^{kN}_{\omega}e(\omega)\|$:

Part 1: Constructions of $S(\omega)$ and $e(\omega)$. These constructions are divided into two steps:

\emph{Step 1}: Construction $S(\omega)$ on $\bigcup_{j=0}^{i-1} \theta^j E_i$ and $e(\omega)$ on $\bigcup_{j=0}^i \theta^j E_i$ for $i=kN+1,\dots$. Let $i\geq kN+1$ and $\omega\in E_i$ be arbitrary.  We write $i=\ell N+r$, where $\ell\in\N$ with $\ell\geq k$ and $0\leq r\leq N-1$. Then, $\bigcup_{j=0}^i \theta^j E_i$ is decomposed as 
\[
\bigcup_{j=0}^i \theta^j E_i
=
\underbrace{\bigcup_{j=0}^{N-1} \theta^j E_i}_{\mathrm{Set\; I}}
\cup 
\underbrace{\bigcup_{j=N}^{\ell N-1} \theta^j E_i}_{\mathrm{Set\; II}}
\cup 
\underbrace{\bigcup_{j=\ell N}^{\ell N+r} \theta^j E_i}_{\mathrm{Set\; III}}.
\]


The constructions of $S(\cdot)$ and $e(\cdot)$ on Set I, Set II, and Set III are given in the following initial step, induction step, and remainder step, respectively.

\noindent 
\underline{Initial step}: Choose and fix an arbitrary random unit map $e:E_i\rightarrow \cH$. Using Lemma \ref{Perturbationalongashortperiod} for a perturbation of $T$ on the disjoint sets $E_i,\theta E_i,\dots,\theta^{N-1}E_i$, there exists $S:\bigcup_{j=0}^{N-1} \theta^j E_i\rightarrow \mathcal K(\cH)$ such that $\|S(\omega)-T(\omega)\|\leq 3\eps (1+\|T\|_{\infty})$ and 
\begin{equation}\label{Construction_Step1a}
\|S_{\omega}^N\|\geq \eps^{N+1}\quad\hbox{and}\quad  \|S_{\omega}^Ne(\omega)\|\geq \eps \|S_{\omega}^N\|\quad \hbox{ for } \omega \in E_i.
\end{equation}
Now, we define $e(\omega)$ on $\theta E_i,\dots,\theta^N E_i$ as 
\begin{equation}\label{Construction_Step1b}
e(\theta^{j}\omega)=\frac{S^{j}_{\omega}e(\omega)}{\|S^j_{\omega}e(\omega)\|}\quad \hbox{ for } \omega \in E_i \hbox{ and } j=1,\dots, N.
\end{equation}

\noindent 
\underline{Induction step}: After the step above, $e(\omega)$ was determined for $\omega\in \theta^N E_i$. Then, applying the construction above to perturb $T$ on the disjoint sets $\theta^{N}E_i,\theta^{N+1}E_i,\dots,\theta^{2N-1}E_i$, $S(\omega)$ is constructed on the sets $\bigcup_{j=N}^{2N-1}\theta^j E_i$ and $e(\omega)$ is constructed on the sets   $\bigcup_{j=N+1}^{2N}\theta^j E_i$ such that \eqref{Construction_Step1a} and \eqref{Construction_Step1b} also hold for $\omega\in \theta^N E_i$. We do this progress inductively and obtain $S$ on $\bigcup_{j=0}^{\ell N-1} \theta^j E_i$ and $e(\omega)$ on $\bigcup_{j=0}^{\ell N} \theta^j E_i$ satisfying that $\|S(\omega)-T(\omega)\|\leq 3\eps(1+ \|T\|_{\infty})$ for  $\omega\in \bigcup_{j=0}^{\ell N-1}\theta^j V_i$ and 
\begin{equation}
\|S_{\omega}^N\|\geq \eps^{N+1}, \|S_{\omega}^Ne(\omega)\|
\geq  \eps \|S_{\omega}^N\|\quad \hbox{for all } \omega \in E_i\cup \theta^N E_i\cup\dots\cup\theta^{\ell N} E_i\label{Nstepnorm},
\end{equation}
and
\begin{equation}\label{Invariance_InductionStep}
e(\theta^j\omega)
=
\frac{S^{j}_{\omega}e(\omega)}{\|S^j_{\omega}e(\omega)\|}\quad \hbox{ for all } \omega \in E_i \hbox{ and } j=1,\dots, \ell N.
\end{equation} 

\noindent 
\underline{Remainder step}: So far, $S(\theta^j\omega)$ has been constructed for $\omega\in E_i$ and $j=0,1,\dots,\ell N-1$ and $e(\theta^j\omega)$ has been constructed for $\omega\in E_i$ and $j=0,1,\dots,\ell N$. For $\omega\in E_i, j=\ell N$, if $\|T(\theta^{j}\omega) e(\theta^{j}\omega)\|\geq \eps$ we let 
\[
S(\theta^{j}\omega)=T(\theta^{j}\omega),\quad   e(\theta^{j+1}\omega)=\frac{S(\theta^j\omega) e(\theta^j\omega)}{\|S(\theta^j\omega) e(\theta^j\omega)\|}
\]
and if $\|T(\theta^{j}\omega) e(\theta^{j}\omega)\|< \eps$ we let
\[
S(\theta^j\omega) v
=
\left\{
\begin{array}{ll}
  (3\eps+T(\theta^j\omega)) v   & \hbox{if } v=e(\theta^j\omega),\\[1ex]
T(\theta^j\omega) v   &  \hbox{if } v\in e(\theta^j\omega)^{\perp},
\end{array}
\right.\quad e(\theta^{j+1}\omega)=\frac{S(\theta^j\omega) e(\theta^j\omega)}{\|S(\theta^j\omega) e(\theta^j\omega)\|}.
\]
Next, we repeat this construction step by step for  $\omega\in E_i, j=\ell N+1,\dots, \ell N+r-1$. Then, $S(\theta^j\omega)$ has been constructed for $\omega \in E_i, j=\ell N,\dots,\ell N+r-1$ and $e(\theta^j\omega)$ has been constructed for $\omega\in E_i, j=\ell N+1,\dots,\ell N+r$. Furthermore, for $j=\ell N,\dots,\ell N+r-1$
\begin{equation}\label{Remainder}
S(\theta^j\omega)e(\theta^j\omega) \hbox{ is colinear to } e(\theta^{j+1}\omega)
\quad \&\quad \|S(\theta^j\omega)e(\theta^j\omega)\|\geq \eps.
\end{equation}
\emph{Step 2}: According to \eqref{PresentationOmega_I} and \eqref{PresentationOmega_II}, it remains to construct $S(\omega)$ on $\bigcup_{j=0}^{N-1}\theta^j U$ and to construct unit random vectors $e(\omega)$ on $\bigcup_{j=1}^{N-1} \theta^j U$. It is noted that $e(\omega)$ was already determined for $\omega\in U\cup\theta^N U$. Thus, by virtue of Proposition \ref{TechnicalProposition_1}(i) we can construct $S(\omega)$ on $\bigcup_{j=0}^{N-1}\theta^j U$ and unit random vectors $e(\omega)$ on  $\bigcup_{j=1}^{N-1} \theta^j U$ such that $\|S(\omega)-T(\omega)\|\leq 3\eps(1+\|T\|_{\infty})$ and 
\begin{equation}\label{Construction_Step2}
S(\theta^i\omega)e(\theta^i\omega)=\beta(\theta^i\omega)e(\theta^{i+1}\omega)\quad \hbox{ for } \omega\in U, i=0,1,\dots,N-1,
\end{equation}
where $\beta: \bigcup_{j=0}^{N-1}\theta^j U\rightarrow \R$ satisfies that $|\beta(\omega)|\geq \eps$ for all $\omega\in \bigcup_{j=0}^{N-1}\theta^j U$.
From \eqref{Construction_Step1b}, \eqref{Invariance_InductionStep}, \eqref{Remainder} and \eqref{Construction_Step2} of  the above construction,  $e(\omega)$ is invariant in the sense that $S(\omega)e(\omega)$ is colinear to $e(\theta\omega)$ for all $\omega\in\Omega$ and $\|S-T\|_{\infty}\leq 3\eps(1+\|T\|_{\infty})$. 

\noindent 
Part 2: Estimation on  $\|S(\omega)e(\omega)\|, \|S^{kN}_{\omega}e(\omega)\|$. Since $\eps\in (0,\frac{1}{6})$ it follows that 
\begin{equation}\label{BoundonS}
\|S\|_{\infty}\leq \|T\|_{\infty}+3\eps(1+\|T\|_{\infty})\leq 1 +2\|T\|_{\infty}.
\end{equation}
We now estimate $\|S(\omega)e(\omega)\|$. For any $\omega\in\bigcup_{j=0}^{N-1}\theta^jU$ we use \eqref{Construction_Step2} to gain that 
\begin{equation}\label{Norm_1}
\|S(\omega)e(\omega)\|
\geq \eps \geq \frac{\eps^{N+1}}{(1+2\|T\|_{\infty})^{N-1}}.
\end{equation}
It remains to estimate $\|S(\omega)e(\omega)\|$ for $\omega\in \bigcup_{i=kN+1}^{\infty}  \Big(E_i\cup \theta E_i\dots\cup\theta^{i-1} E_i\Big)$. Let $\omega=\theta^j \widehat\omega$ for some $\widehat\omega \in E_i$, where $i=\ell N+r \geq kN+1$ and $0\leq j\leq i-1$. Write $j=sN+t$. In the case that $s<\ell$,  by the construction in the Initial Step and Induction Step of Part 1, we have  
\begin{align*}
\|S(\omega)e(\omega)\|
&=\|S(\theta^{sN+t}\widehat\omega)e(\theta^{sN+t}\widehat\omega)\|\notag\\
&=\frac{\|S^N_{\theta^{sN}\widehat\omega}e(\theta^{sN}\widehat\omega)\|}{\|S^t_{\theta^{sN}\widehat\omega}e(\theta^{sN}\widehat\omega)\|\|S^{N-t-1}_{\theta^{sN+t+1}\widehat\omega}e(\theta^{sN+t+1}\widehat\omega)\|}\notag\\
&\geq  
\frac{\eps^{N+1}}{\|S\|_{\infty}^{N-1}},
\end{align*}
which together with \eqref{Norm_1} implies that
\begin{equation}\label{Norm_2}
\|S(\omega)e(\omega)\| 
\geq \frac{\eps^{N+1}}{(1+2\|T\|_{\infty})^{N-1}}.
\end{equation}
In the case that $s=\ell$, by the construction in the Remainder Step, we have 
\begin{equation}\label{Norm_3}
\|S(\omega)e(\omega)\|
\geq \eps \geq \frac{\eps^{N+1}}{(1+2\|T\|_{\infty})^{N-1}}.
\end{equation}
Combining \eqref{Norm_1}, \eqref{Norm_2} and \eqref{Norm_3}, the estimate on $\|S(\omega)e(\omega)\|$ as in  \eqref{Lowerestimate} was verified. To conclude the proof, we show that $\|S^{kN}_{\omega}e(\omega)\|$ fulfills  \eqref{Goodgrowthrate}. Thanks to \eqref{PresentationOmega_II}, we can verify \eqref{Goodgrowthrate} by considering  two separated cases $\omega\in  \bigcup_{j=0}^{N-1}\theta^j U$ (Case 1) or $\omega\in \bigcup_{i=kN+1}^{\infty}  \Big(E_i\cup \theta E_i\dots\cup\theta^{i-1} E_i\Big)$ (Case 2).

\noindent
\emph{Case 1}: $\omega\in \bigcup_{j=0}^{N-1}\theta^j U$. Let $\omega=\theta^{\tau}\widehat \omega$ for some $\widehat\omega\in U$ and $\tau\in\{0,\dots,N-1\}$. Then, by \eqref{Construction_Step2} we have 
\begin{equation}\label{Case1_Eq1}
\|S_{\omega}^{kN} e(\omega)\|=\|S_{\theta^{\tau}\widehat\omega}^{kN}e(\theta^{\tau}\widehat\omega)\|\geq \eps^{N-\tau} \|S_{\theta^N\widehat\omega}^{kN-(N-\tau)}e(\theta^N\widehat\omega)\|.
\end{equation}
Note that $\theta^N\widehat\omega\in \theta^N U$ and from the invariance of $e(\omega)$, \eqref{Nstepnorm} and \eqref{Nstepnorm}  we derive that 
\begin{align}
&\|S_{\theta^N\widehat\omega}^{(k-1)N+\tau}e(\theta^N\widehat\omega)\|\notag\\[1ex]
=&\|S_{\theta^{kN}\widehat\omega}^{\tau}e(\theta^{kN}\widehat\omega)\| \|S^{N}_{\theta^{(k-1)N}\widehat\omega}e(\theta^{(k-1)N}\widehat\omega)\|\dots\|S^N_{\theta^N\widehat\omega}e(\theta^N\widehat\omega)\|\notag\\[1ex]
\geq& 
\eps^{k-1}\|S_{\theta^{kN}\widehat\omega}^{\tau}e(\theta^{kN}\widehat\omega)\| \|S^{N}_{\theta^{(k-1)N}\widehat\omega}\|\dots\|S^N_{\theta^{2N}\widehat\omega}\|\|S^N_{\theta^N\widehat\omega}\|.\label{FirstEstimate}
\end{align}
By \eqref{Nstepnorm} we have 
\[
\|S_{\theta^{kN}\widehat\omega}^{\tau}e(\theta^{kN}\widehat\omega)\|\geq \frac{\eps^{N+1}}{\|S_{\theta^{kN+\tau}\widehat\omega}^{N-\tau}\|}
\geq 
\frac{\eps^{N+1}}{\|S\|_{\infty}^{N-\tau}},
\]
which together with the fact that $\|S_{\theta^{kN}\widehat\omega}^{\tau}\|\leq \|S\|_{\infty}^{\tau}$ implies that 
\[
\|S_{\theta^{kN}\widehat\omega}^{\tau}e(\theta^{kN}\widehat\omega)\|
\geq
\frac{\eps^{N+1}}{\|S\|_{\infty}^{N}} \||S_{\theta^{kN}\widehat\omega}^{\tau}\|.
\]
Combining the above inequality with \eqref{Case1_Eq1} and \eqref{FirstEstimate} yields that 
\begin{align*}
\|S_{\omega}^{kN}e(\omega)\|
\geq &
\eps^{N-\tau}\eps^{k-1}
\frac{\eps^{N+1}}{\|S\|_{\infty}^N} \||S_{\theta^{kN}\widehat\omega}^{\tau}\|\|S^{N}_{\theta^{(k-1)N}\widehat\omega}\|\dots\|S^N_{\theta^{2N}\widehat\omega}\|\|S^N_{\theta^N\widehat\omega}\|\\[1ex]
\geq &
\frac{\eps^{2N+k-\tau}}{\|S\|_{\infty}^{N-\tau}}\|S_{\theta^N\widehat\omega}^{(k-1)N+\tau}\|.
\end{align*}
Note that $\|S_{\omega}^{kN}\|=\|S_{\theta^N\widehat\omega}^{(k-1)N+\tau}S_{\theta^{\tau}\widehat\omega}^{N-\tau}\|
\leq \|S_{\theta^N\widehat\omega}^{(k-1)N+\tau}\|\|S\|_{\infty}^{N-\tau}$. Thus, we arrive at 
\[
\|S_{\omega}^{kN}e(\omega)\|
\geq 
\frac{\eps^{2N+k-\tau}}{\|S\|_{\infty}^{2(N-\tau)}}\|S_{\omega}^{kN}\|
\geq 
\frac{\eps^{2N+k}}{( 1+2\|T\|_{\infty})^{2N}}\|S_{\omega}^{kN}\|.
\]
It means that \eqref{Goodgrowthrate} holds for all $\omega\in \bigcup_{j=0}^{N-1}\theta^j U$.

\noindent
\emph{Case 2}: $\omega\in \bigcup_{i=kN+1}^{\infty}  \Big(E_i\cup \theta E_i\dots\cup\theta^{i-1} E_i\Big)$. Let $i\geq kN+1$ and $\tau\in \{0,\dots,i-1\}$ such that $\omega\in\theta^{\tau }E_i$. We write $\omega=\theta^{\tau}\widehat\omega$ for some $\widehat\omega\in E_i$ and let  $i=\ell N+r, \tau=sN+t$, where $s,\ell\in\N$ and $r,t\in \{0,\dots,N-1\}$. To estimate $\|S_{\omega}^{kN}e(\omega)\|=
\|S^{kN}_{\theta^{sN+t}\widehat\omega}e(\theta^{sN+t}\widehat\omega)\|$, we consider three subcases $\tau+kN\leq i-1$ (Subcase 2a), $i\leq \tau+kN\leq i+N-1$ (Subcase 2b) and $\tau+kN\geq i+N$ (Subcase 2c).

\emph{Subcase 2a}: $\tau+kN\leq i-1$. Then, by invariance of $e(\omega)$ we have 
\begin{align*}
&\|S^{kN}_{\theta^{sN+t}\widehat\omega}e(\theta^{sN+t}\widehat\omega)\|\\[1ex]
=&\|S^{N-t}_{\theta^{sN+t}\widehat\omega}e(\theta^{sN+t}\widehat\omega)\|\|S^{(k-1)N}_{\theta^{(s+1)N}\widehat\omega}e(\theta^{(s+1)N}\widehat\omega)\|
\|S^{t}_{\theta^{(k+s)N}\widehat\omega}e(\theta^{(k+s)N}\widehat\omega)\|,
\end{align*}
which together with \eqref{Construction_Step1a} and \eqref{Nstepnorm} implies that 
\begin{align}
&\|S^{kN}_{\theta^{sN+t}\widehat\omega}e(\theta^{sN+t}\widehat\omega)\|\notag
\\[1ex]
\geq& 
\eps^{k-1} \|S^{N-t}_{\theta^{sN+t}\widehat\omega}e(\theta^{sN+t}\widehat\omega)\|\|S^{(k-1)N}_{\theta^{(s+1)N}\widehat\omega}\|
\|S^{t}_{\theta^{(k+s)N}\widehat\omega}e(\theta^{(k+s)N}\widehat\omega)\|.
\label{Case2a_Eq1}
\end{align}
Using again  \eqref{Construction_Step1a} and \eqref{Nstepnorm}, we obtain that 
\begin{equation}\label{Subcase2a:Eq2}
\|S^{N-t}_{\theta^{sN+t}\widehat\omega}e(\theta^{sN+t}\widehat\omega)\|=\frac{\|S^{N}_{\theta^{sN}\widehat\omega}e(\theta^{sN}\widehat\omega)\|}{\|S^{t}_{\theta^{sN}\widehat\omega}e(\theta^{sN}\widehat\omega)\|}
\geq 
\frac{\eps^{N+1}}{\|S\|_{\infty}^t}.
\end{equation}
It remains to estimate $\|S^{t}_{\theta^{(k+s)N}\widehat\omega}e(\theta^{(k+s)N}\widehat\omega)\|$. Since $\tau+k\ell\leq i-1$ it follows that  either $k+s<\ell$ or $k+s=\ell$. Concerning the first case, i.e. $k+s<\ell$, from  \eqref{Construction_Step1a} and \eqref{Nstepnorm} we derive that 
\[
\|S^{t}_{\theta^{(k+s)N}\widehat\omega}e(\theta^{(k+s)N}\widehat\omega)\|
\geq 
\frac{\|S^{N}_{\theta^{(k+s)N}\widehat\omega}e(\theta^{(k+s)N}\widehat\omega)\|}{\|S^{N-t}_{\theta^{(k+s)N}\widehat\omega}e(\theta^{(k+s)N}\widehat\omega)\|}
\geq\frac{\eps^{N+1}}{\|S\|_{\infty}^{N-t}},
\]
which together with \eqref{Case2a_Eq1} and \eqref{Subcase2a:Eq2} implies that 
\begin{equation}\label{Firstcase}
\|S^{kN}_{\theta^{sN+t}\widehat\omega}e(\theta^{sN+t}\widehat\omega)\|
\geq \frac{\eps^{2N+k+1}}{\|S\|_{\infty}^N}\|S^{kN}_{\omega}\|
\geq 
\frac{\eps^{2N+k+1}}{(1+2\|T\|_{\infty})^N}\|S^{kN}_{\omega}\|,
\end{equation}
where we used \eqref{BoundonS} to obtain the last inequality. Meanwhile, in the other case, i.e. $k+s=\ell$, by using \eqref{Remainder} we obtain that 
\[
\|S^{t}_{\theta^{(k+s)N}\widehat\omega}e(\theta^{(k+s)N}\widehat\omega)\|
\geq 
\eps^t
\geq 
\eps^t\frac{\|S^{t}_{\theta^{(k+s)N}\widehat\omega}\|}{\|S\|_{\infty}^{t}},
\]
which together with \eqref{Case2a_Eq1} and \eqref{Subcase2a:Eq2} implies that 
\begin{equation}\label{Secondcase}
\|S^{kN}_{\theta^{sN+t}\widehat\omega}e(\theta^{sN+t}\widehat\omega)\|
\geq \frac{\eps^{N+k+t}}{\|S\|_{\infty}^N}\|S^{kN}_{\omega}\|
\geq 
\frac{\eps^{2N+k}}{(1+2\|T\|_{\infty})^N}\|S^{kN}_{\omega}\|,
\end{equation}
where we used \eqref{BoundonS} to obtain the last inequality. Combining \eqref{Firstcase} and \eqref{Secondcase}, we obtain that in this subcase 
\[
\|S^{kN}_{\theta^{sN+t}\widehat\omega}e(\theta^{sN+t}\widehat\omega)\|
\geq 
\frac{\eps^{2N+k+1}}{(1+2\|T\|_{\infty})^N}\|S^{kN}_{\omega}\|,
\]
which means that  \eqref{Goodgrowthrate} holds. 

\emph{Subcase 2b}: $i\leq \tau+kN\leq i+N-1$. Hence, from $i=\ell N+r, \tau=sN+t$ with $r,t\in \{0,\dots,N-1\}$ we have 
\begin{equation}\label{estimateindex}
0\leq kN+\tau-i\leq N-1,\quad (\ell-s)\leq k    
\end{equation}
By invariance of $e(\omega)$ we have 
\begin{align*}
\|S^{kN}_{\theta^{sN+t}\widehat\omega}e(\theta^{sN+t}\widehat\omega)\|
&=\|S^{N-t}_{\theta^{sN+t}\widehat\omega}e(\theta^{sN+t}\widehat\omega)\|\|S^{(\ell-s-1)N}_{\theta^{(s+1)N}\widehat\omega}e(\theta^{(s+1)N}\widehat\omega)\|\\[1ex]
&\times
\|S^{r}_{\theta^{\ell N}\widehat\omega}e(\theta^{\ell N}\widehat\omega)\|\|S^{kN+\tau-i}_{\theta^i\widehat\omega}e(\theta^i\widehat\omega)\|.
\end{align*}
Using \eqref{Construction_Step1a} and  \eqref{Nstepnorm}, we obtain that 
\begin{align}
\|S^{N-t}_{\theta^{sN+t}\widehat\omega}e(\theta^{sN+t}\widehat\omega)\|
&= \frac{\|S^{N}_{\theta^{sN}\widehat\omega}e(\theta^{sN}\widehat\omega)\|}{\|S^{t}_{\theta^{sN}\widehat\omega}e(\theta^{sN}\widehat\omega)\|}\notag\\
&\geq \frac{\eps^{N+1}}{\|S\|_{\infty}^N}\|S^{N-t}_{\theta^{sN+t}\widehat\omega}\|,\label{Subcase2b_Eq1}\\[1ex]
\|S^{(\ell-s-1)N}_{\theta^{(s+1)N}\widehat\omega}e(\theta^{(s+1)N}\widehat\omega)\|
&=\prod_{j=s+1}^{\ell-1}\|S_{\theta^{jN}\widehat\omega}e(\theta^{jN}\widehat\omega)\|\notag\\
&\geq 
\eps^{\ell-s-1} \|S^{(\ell-s-1)N}_{\theta^{(s+1)N}\widehat\omega}\|\label{Subcase2b_Eq2}.
\end{align}
Meanwhile, by \eqref{Remainder} we have 
\begin{equation}\label{Subcase2b_Eq3}
\|S^{r}_{\theta^{\ell N}\widehat\omega}e(\theta^{\ell N}\widehat\omega)\|
\geq \eps^r\geq \eps^r \frac{\|S^{r}_{\theta^{\ell N}\widehat\omega}\|}{\|S\|_{\infty}^r}.
\end{equation}
Finally, note that $\theta^i\widehat\omega\in U$ and thus by using \eqref{Construction_Step2} we gain that 
\begin{equation}\label{Subcase2b_Eq4}
\|S^{kN+\tau-i}_{\theta^i\widehat\omega}e(\theta^i\widehat\omega)\|\geq \eps^{kN+\tau-i}
\geq \eps^{kN+\tau-i}\frac{\|S^{kN+\tau-i}_{\theta^i\widehat\omega}\|}{\|S\|_{\infty}^{kN+\tau-i}}.
\end{equation}
Combining \eqref{Subcase2b_Eq1}, \eqref{Subcase2b_Eq2}, \eqref{Subcase2b_Eq3} and \eqref{Subcase2b_Eq4} yields that \[
\|S^{kN}_{\theta^{sN+t}\widehat\omega}e(\theta^{sN+t}\widehat\omega)\|
\geq 
\frac{\eps^{N+(\ell-s)+r+(kN+\tau-i)}}{\|S\|_{\infty}^{N+r+(kN+\tau-i)}}\|S^{kN}_{\theta^{sN+t}\widehat\omega}\|,
\]
which together with \eqref{estimateindex} and \eqref{BoundonS} implies that 
\[
\|S^{kN}_{\theta^{sN+t}\widehat\omega}e(\theta^{sN+t}\widehat\omega)\|
\geq 
\frac{\eps^{3N+k}}{(1+2\|T\|_{\infty})^{3N}}\|S^{kN}_{\theta^{sN+t}\widehat\omega}\|.
\]
Thus, \eqref{Goodgrowthrate} holds in this subcase. 

\emph{Subcase 2c}: $ \tau+kN\geq i+N$.  Hence, from $i=\ell N+r, \tau=sN+t$ with $r,t\in \{0,\dots,N-1\}$ we have 
\begin{equation}\label{estimateindex_casec}
0\leq (\ell-s)\leq k-1.    
\end{equation}
By invariance of $e(\omega)$ we have 
\begin{equation}\label{Twoparts}
\|S^{kN}_{\theta^{sN+t}\widehat\omega}e(\theta^{sN+t}\widehat\omega)\|=  \|S^{i-(sN+t)}_{\theta^{sN+t}\widehat\omega}e(\theta^{sN+t}\widehat\omega)\|  \|S^{kN+(sN+t)-i}_{\theta^{i}\widehat\omega}e(\theta^{i}\widehat\omega)\|.
\end{equation}
On the one hand, analogous to the estimates in \eqref{Subcase2b_Eq1}, \eqref{Subcase2b_Eq2} and \eqref{Subcase2b_Eq3} we have 
\begin{equation}\label{Estimate_Part1}
  \|S^{i-(sN+t)}_{\theta^{sN+t}\widehat\omega}e(\theta^{sN+t}\widehat\omega)\| 
  \geq \frac{\eps^{N+r+(\ell-s)}}{\|S\|_{\infty}^{N+r}}\|S^{i-(sN+t)}_{\theta^{sN+t}\widehat\omega}\|.
\end{equation}
On the other hand,  by denoting $\widetilde\omega:=\theta^i\widehat\omega$ we have $\widetilde \omega\in U$ and from  invariance of $e(\omega)$ we derive that 
\begin{align}
\|S^{kN+(sN+t)-i}_{\theta^{i}\widehat\omega}e(\theta^{i}\widehat\omega)\|
&=\|S_{\widetilde\omega}^{(k+s-\ell)N+(t-r)}e(\widetilde\omega)\|\notag\\[1ex]
&=
\|S_{\theta^N\widetilde\omega}^{(k+s-\ell-1)N+(t-r)}e(\theta^N\widetilde\omega)\|
\|S_{\widetilde\omega}^{N}e(\widetilde\omega)\|.\label{New_Estimate1}
\end{align}
Thanks to the construction of $S(\omega)$ for $\omega\in\bigcup_{j=0}^{N-1}\theta^j U$ as in Step 2, we have \begin{equation}\label{New_Estimate2}
 \|S_{\widetilde\omega}^{N}e(\widetilde\omega)\|
 \geq 
 \eps^N\geq \frac{\eps^N}{\|S\|_{\infty}^N}\|S^N_{\widetilde\omega}\|.
\end{equation}
Thus, to estimate $\|S^{kN+(sN+t)-i}_{\theta^{i}\widehat\omega}e(\theta^{i}\widehat\omega)\|$ it remains to estimate $\|S_{\theta^N\widetilde\omega}^{(k+s-\ell-1)N+(t-r)}e(\theta^N\widetilde\omega)\|$. In fact, we show that 
\begin{equation}\label{Smallaim}
  \|S_{\theta^N\widetilde\omega}^{(k+s-\ell-1)N+(t-r)}e(\theta^N\widetilde\omega)\| 
  \geq \frac{\eps^{k+s-\ell+N}}{\|S\|_{\infty}^N}       \|S_{\theta^N\widetilde\omega}^{(k+s-\ell-1)N+(t-r)}\|.
\end{equation}
For this purpose, we consider two separated cases $(t-r)\geq 0$ or $(t-r)<0$:
\begin{itemize}
    \item If $(t-r)\geq 0$ then
    \begin{align*}
       &\|S_{\theta^N\widetilde\omega}^{(k+s-\ell-1)N+(t-r)}e(\theta^N\widetilde\omega)\|\\[1ex]
       &= \|S_{\theta^N\widetilde\omega}^{(k+s-\ell-1)N}e(\theta^N\widetilde\omega)\|
       \frac{\|S^{N}_{\theta^{(k+s-\ell)N}\widetilde\omega}e(\theta^{(k+s-\ell)N}\widetilde\omega)\|}{\|S^{N-(t-r)}_{\theta^{(k+s-\ell)N+(t-r)}\widetilde\omega}e(\theta^{(k+s-\ell)N+(t-r)}\widetilde\omega)\|},
    \end{align*}
which together with \eqref{Construction_Step1a} and \eqref{Nstepnorm} implies that  \begin{align*}
       \|S_{\theta^N\widetilde\omega}^{(k+s-\ell-1)N+(t-r)}e(\theta^N\widetilde\omega)\|
       &\geq \eps^{k+s-\ell-1}\|S_{\theta^N\widetilde\omega}^{(k+s-\ell-1)N}\|\frac{\eps^{N+1}}{\|S\|_{\infty}^{N-(t-r)}},\\[1ex]
       &\geq 
       \frac{\eps^{k+s-\ell+N}}{\|S\|_{\infty}^N}       \|S_{\theta^N\widetilde\omega}^{(k+s-\ell-1)N+(t-r)}\|.
    \end{align*}
Thus, \eqref{Smallaim} is proved in this case.
\item If $(t-r)<0$ then 
 \begin{align*}
       &\|S_{\theta^N\widetilde\omega}^{(k+s-\ell-1)N+(t-r)}e(\theta^N\widetilde\omega)\|\\[1ex]
       &= \|S_{\theta^N\widetilde\omega}^{(k+s-\ell-2)N}e(\theta^N\widetilde\omega)\|
       \frac{\|S^{N}_{\theta^{(k+s-\ell-1)N}\widetilde\omega}e(\theta^{(k+s-\ell)N}\widetilde\omega)\|}{\|S^{r-t}_{\theta^{(k+s-\ell)N+(t-r)}\widetilde\omega}e(\theta^{(k+s-\ell)N+(t-r)}\widetilde\omega)\|},
    \end{align*}
which together with \eqref{Construction_Step1a} and \eqref{Nstepnorm} implies that  \begin{align*}
       \|S_{\theta^N\widetilde\omega}^{(k+s-\ell-1)N+(t-r)}e(\theta^N\widetilde\omega)\|
       &\geq \eps^{k+s-\ell-2}\|S_{\theta^N\widetilde\omega}^{(k+s-\ell-2)N}\|\frac{\eps^{N+1}}{\|S\|_{\infty}^{r-t}},\\[1ex]
       &\geq 
       \frac{\eps^{k+s-\ell+N-1}}{\|S\|_{\infty}^N}       \|S_{\theta^N\widetilde\omega}^{(k+s-\ell-1)N+N-(r-t)}\|\\
       &\geq \frac{\eps^{k+s-\ell+N}}{\|S\|_{\infty}^N}       \|S_{\theta^N\widetilde\omega}^{(k+s-\ell-1)N+(t-r)}\|
    \end{align*}
Thus, \eqref{Smallaim} is proved.
\end{itemize}
Combining \eqref{Smallaim} with \eqref{New_Estimate1} and \eqref{New_Estimate2}, we obtain that 
\[
\|S^{kN+(sN+t)-i}_{\theta^{i}\widehat\omega}e(\theta^{i}\widehat\omega)\|
\geq 
       \frac{\eps^{k+s-\ell+2N}}{\|S\|_{\infty}^{2N}} \|S^{kN+(sN+t)-i}_{\theta^{i}\widehat\omega}\|,
\]
which together with \eqref{Twoparts} and \eqref{Estimate_Part1} leads to
\begin{align*}
\|S^{kN}_{\theta^{sN+t}\widehat\omega}e(\theta^{sN+t}\widehat\omega)\|
& \geq \frac{\eps^{k+3N+r}}{\|S\|_{\infty}^{3N+r}}\|S^{kN}_{\theta^{sN+t}\widehat\omega}\|\\[1ex]
&\geq 
\frac{\eps^{k+4N}}{(1+2\|T\|_{\infty})^{4N}} \|S^{kN}_{\theta^{sN+t}\widehat\omega}\|,
\end{align*}
where we used \eqref{BoundonS} and the fact that $\eps\in (0,\frac{1}{6})$ and $r<N$ to obtain the last inequality. It means that \eqref{Goodgrowthrate} also holds in this case, and the proof is complete.
\end{proof}

\end{document}
